\documentclass[10pt,a4paper]{amsart}
\usepackage[margin=19mm]{geometry}
\usepackage{amsmath,amssymb,mathtools}
\usepackage[scr=rsfs,cal=euler]{mathalfa}
\usepackage{xfrac}
\usepackage[unicode,hidelinks,bookmarks=false]{hyperref}
\newcommand{\ie}{i.e.\ }
\newcommand{\cf}{cf.\ }
\newcommand{\resp}{respectively\ }
\newcommand{\Subsection}{\subsection}
\providecommand{\nobreakdash}{\nobreak\hbox{-}}
\setlength{\emergencystretch}{2em}
\multlinegap0pt
\usepackage{amssymb}
\usepackage{mathtools}
\newtheorem{proposition}{Proposition}
\newcommand{\K}{\mathbb{K}}
\renewcommand{\ge}{\geqslant}
\renewcommand{\le}{\leqslant}
\title{Raja's covering index of $L_p$ spaces}
\author{Tomasz Kania, Natalia Ma\'slany}
\date{Source: arXiv:2512.13249v2 (CC BY 4.0)}
\begin{document}
\maketitle

\noindent\emph{Source and licence.} Raja's covering index of $L_p$ spaces. Version arXiv:2512.13249v2 (CC BY 4.0), distributed by arXiv under the Creative Commons Attribution 4.0 International licence, http://creativecommons.org/licenses/by/4.0/, as recorded in the arXiv licence metadata for this exact version and shown on the abstract page https://arxiv.org/abs/2512.13249v2. Full licence notice: \texttt{licenses/arxiv-2512.13249-CC-BY-4.0.txt}.\par\smallskip

\noindent\textit{Editorial scope.} Counted unit: the article's proposition prop:Lp-mu-upper (source0.tex lines 501--508). The article's complete original proof of it is delivered with it (source0.tex lines 510--695, one \texttt{proof} environment of 186 lines). No mathematics is written, repaired or re-derived: the statement and the proof are the article's own lines, in the article's own notation, and no retained line is substituted or re-flowed.  No further source-local context is needed: every cross-reference of every retained line resolves inside the two delivered spans.  Nothing was omitted from any retained span, so the package carries no omission record.  No retained line contains a citation, so the wrapper carries no bibliography and no bibliography entry.  No retained line contains a figure environment, an image-inclusion command or drawing code, so no image is delivered and the package declares no asset dependency.  Editorial formatting only, in addition: an AMS \texttt{amsart} preamble that restates the article's own declarations and macros used by the retained lines, namely the environments \texttt{proposition} on separate counters, together with the standard packages those lines need.  The retained environments print in this document as Proposition 1 in order of appearance, whereas the article prints the same items with the numbering of its own full document.  The complete native source of the article as distributed in the arXiv e-print for this exact version is bundled unchanged as \texttt{sources/191012/source0.tex}.
\begin{proposition}\label{prop:Lp-mu-upper}
Let $1\le p<\infty$ and let $(\Omega,\Sigma,\mu)$ be a $\sigma$-finite
measure space such that $L_p(\mu)$ is infinite-dimensional. Then, for every
$n\in\mathbb{N}$,
\[
\Theta_{L_p(\mu)}(n)\ \le\ n^{-1/p}.
\]
\end{proposition}

\begin{proof}
Fix $n\in\mathbb N$, and choose a partition
\[
\Omega=\bigsqcup_{k=1}^n E_k
\]
such that $L_p(E_k)$ is infinite-dimensional for every $k$.

For each $k=1,\dots,n$ define the projection
\[
P_k f := f\cdot\mathbf{1}_{E_k}, \qquad f\in L_p(\mu).
\]
Then $P_k:L_p(\mu)\to L_p(\mu)$ is a contractive projection.

For $f\in L_p(\mu)$ we have
\[
\sum_{k=1}^n \|P_k f\|_p^p
= \sum_{k=1}^n \int_{E_k} |f|^p\, d\mu
= \int_{\Omega} |f|^p\, d\mu
= \|f\|_p^p.
\]
In particular, if $f\in B_{L_p(\mu)}$, then
\[
\sum_{k=1}^n \|P_k f\|_p^p \le 1.
\]

Define, for $1\le k\le n$,
\[
A_k := \bigl\{ f\in B_{L_p(\mu)}:\ \|P_k f\|_p\le n^{-1/p} \bigr\}.
\]
Each $A_k$ is closed, convex and bounded.

If $f\in B_{L_p(\mu)}$, then $\|P_k f\|_p^p \le 1/n$ for some $k$, that
is, $\|P_k f\|_p\le n^{-1/p}$, and hence $f\in A_k$. Thus
\[
B_{L_p(\mu)} = \bigcup_{k=1}^n A_k.
\]

\smallskip
\noindent\emph{Lower bound on $\varrho(A_k)$.}
If $\|g\|_p\le n^{-1/p}$, then $\|P_k g\|_p\le \|g\|_p\le n^{-1/p}$, so
$g\in A_k$. Taking $x=0$ and $Y=L_p(\mu)$ in the definition of the
essential inradius, we get
\[
n^{-1/p} B_{L_p(\mu)} \subset A_k,
\]
so $\varrho(A_k)\ge n^{-1/p}$.

\smallskip
\noindent\emph{Upper bound on $\varrho(A_k)$.}
Suppose, for a contradiction, that $\varrho(A_k)>n^{-1/p}$. Then there exist
$f\in A_k$, a finite-codimensional subspace $Y\le L_p(\mu)$, and
$r>n^{-1/p}$ such that
\[
f + r(B_{L_p(\mu)}\cap Y)\subset A_k.
\]
Let $v := P_k f\in L_p(E_k)$; then $\|v\|_p\le n^{-1/p}$.

Since $Y$ has finite codimension, there exist $\varphi_1,\dots,\varphi_m\in
L_p(\mu)^*$ such that
\[
Y = \bigcap_{i=1}^m \ker\varphi_i.
\]

Since $L_p(E_k)$ is infinite-dimensional and $(E_k,\Sigma|_{E_k},\mu|_{E_k})$
is $\sigma$-finite, there exists a sequence $(C_q)_{q\ge1}$ of pairwise
disjoint measurable subsets of $E_k$ such that
\[
0<\mu(C_q)<\infty\qquad(q\ge1).
\]
For $j\ge1$ and $0\le \ell\le m$, set
\[
B_{j,\ell}:=C_{(m+1)(j-1)+\ell+1},
\qquad
U_j:=\bigcup_{\ell=0}^m B_{j,\ell}.
\]
Then the sets $U_j$ are pairwise disjoint. Set
\[
a_j := \|v\mathbf 1_{U_j}\|_p.
\]
Since the $U_j$ are disjoint and $v\in L_p(E_k)$, we have
\[
\sum_{j=1}^\infty a_j^p
=
\sum_{j=1}^\infty \int_{U_j}|v|^p\,d\mu
\le \|v\|_p^p<\infty,
\]
hence $a_j\to0$.

For each $j$, define
\[
f_{j,\ell}:=\mathbf 1_{B_{j,\ell}},\qquad \ell=0,\dots,m,
\]
and consider the linear map
\[
T_j:\K^{m+1}\longrightarrow\K^m,\qquad
T_j(\alpha_0,\dots,\alpha_m)
=
\Bigl(\varphi_i\Bigl(\sum_{\ell=0}^m \alpha_\ell f_{j,\ell}\Bigr)\Bigr)_{i=1}^m.
\]
Since $\dim \K^{m+1}>\dim \K^m$, the kernel of $T_j$ is non-trivial, so
there exists a non-zero
\[
w_j=\sum_{\ell=0}^m \alpha_\ell f_{j,\ell}\in Y
\]
supported in $U_j$. After normalising, we may assume that
\[
u_j:=\frac{w_j}{\|w_j\|_p}
\]
satisfies $\|u_j\|_p=1$ and is supported in $U_j$.

Fix $r>n^{-1/p}$ as above and consider the continuous function
\[
F(a)
=
\|v\|_p^p - a^p + (r-a)^p,\qquad a\ge0.
\]
We have
\[
F(0)=\|v\|_p^p+r^p>r^p>1/n,
\]
so, by continuity, there exists $\delta>0$ such that
\[
0<\delta<r
\quad\text{and}\quad
F(a)>1/n \ \text{whenever}\ 0\le a\le\delta.
\]

Choose $j_0$ large enough so that $a_{j_0}\le\delta$, and set
\[
u:=u_{j_0},\qquad U:=U_{j_0}.
\]
Then $u\in Y$, $\|u\|_p=1$, and $\|v\mathbf 1_U\|_p=a_{j_0}\le\delta$.

Decompose
\[
v=v_1+v_2,
\qquad
v_1=v\mathbf 1_{E_k\setminus U},\quad
v_2=v\mathbf 1_U.
\]
Since $u$ is supported in $U$, we have
\[
\|v+ru\|_p^p
=
\|v_1\|_p^p+\|v_2+ru\|_p^p.
\]
By the triangle inequality in $L_p$,
\[
\|v_2+ru\|_p
\ge r-\|v_2\|_p
= r-a_{j_0},
\]
and therefore
\[
\|v+ru\|_p^p
\ge
\|v\|_p^p-a_{j_0}^p+(r-a_{j_0})^p
=
F(a_{j_0})
>
1/n.
\]

Finally, set $y:=ru\in Y$. Then $\|y\|_p=r$ and
\[
\|P_k(f+y)\|_p^p
=
\|v+ru\|_p^p
>
1/n,
\]
so $\|P_k(f+y)\|_p>n^{-1/p}$, and therefore $f+y\notin A_k$. This
contradicts the assumption that
\[
f+r(B_{L_p(\mu)}\cap Y)\subset A_k.
\]

Thus no such $r>n^{-1/p}$ exists and $\varrho(A_k)\le n^{-1/p}$. Combining
with the lower bound we obtain $\varrho(A_k)=n^{-1/p}$ for each $k$, and
hence
\[
\Theta_{L_p(\mu)}(n)
\le \max_{1\le k\le n}\varrho(A_k)
= n^{-1/p}.
\]
\end{proof}

\end{document}
